% !PRINTED 368
\weberpage{368}
By what was remarked above, however, there are also always solutions of the equation (4) when $\gamma$ is divisible by the greatest common divisor of $\alpha$ and $\beta$. In particular we can state the theorem: \gesperrt{The greatest common divisor $\delta$ of two numbers $\alpha$, $\beta$ can be represented in the form}
\[
\alpha x - \beta y = \delta
\]
\gesperrt{in which $x$, $y$ are integers.}

Theorem III can be generalised in the following way:

\gesperrt{If $\alpha_1$, $\alpha_2$, $\alpha_3 \ldots$ are given integers in arbitrary number, without a divisor common to them all, and $m$ is an arbitrarily given number, then the integers $x_1$, $x_2$, $x_3 \ldots$ can always be determined so that}

\weq{5}{m = \alpha_1 x_1 + \alpha_2 x_2 + \alpha_3 x_3 + \cdots}
\gesperrt{holds.}

If the right-hand side of (5) consists of only two terms, this theorem coincides with Theorem III. If, then, we assume the possibility of satisfying a formula (5) as proved whenever the right-hand side consists of fewer terms, we can, if $\delta$ is the greatest common divisor of $\alpha_2$, $\alpha_3 \ldots$, since $\delta$ must be relatively prime to $\alpha_1$, satisfy the two equations:
\[
\delta = \alpha_2 y_2 + \alpha_3 y_3 + \cdots
\]
\[
m = \alpha_1 x_1 + \delta x
\]
by integral $x$, $x_1$, $y_2$, $y_3 \ldots$ But then the equation (5) is satisfied by
\[
x_1, \quad x_2 = x y_2, \quad x_3 = x y_3 \ldots
\]

The problem solved by Theorem III is also expressed in number theory as follows: If $\alpha$, $\beta$, $\gamma$ are given numbers, a number $y$ is to be found which satisfies the congruence
\[
\alpha y \equiv \gamma \ (\mathrm{mod}\ \beta)
\]
This problem therefore always has a solution when $\gamma$ is divisible by the greatest common divisor of $\alpha$ and $\beta$.
